\documentclass[aspectratio=169,11pt]{beamer}

% ============================================================
% Theme and typography
% ============================================================

\usetheme{Madrid}
\usecolortheme{default}
\usefonttheme{professionalfonts}

\setbeamertemplate{navigation symbols}{}
\setbeamertemplate{headline}{}

\setbeamertemplate{footline}{%
  \leavevmode\hbox{%
  \begin{beamercolorbox}
    [wd=.84\paperwidth,ht=2.7ex,dp=1.1ex,leftskip=.8em]
    {author in head/foot}
    CSCI 6461 --- Computer Systems Architecture
  \end{beamercolorbox}%
  \begin{beamercolorbox}
    [wd=.16\paperwidth,ht=2.7ex,dp=1.1ex,rightskip=.8em]
    {date in head/foot}
    \hfill\insertframenumber/\inserttotalframenumber
  \end{beamercolorbox}}
}

\setbeamersize{
  text margin left=0.65cm,
  text margin right=0.65cm
}

\setbeamerfont{frametitle}{
  size=\Large,
  series=\bfseries
}

\setbeamerfont{block title}{
  series=\bfseries
}

\setbeamertemplate{blocks}[rounded][shadow=false]

% ============================================================
% Packages
% ============================================================

\usepackage[T1]{fontenc}
\usepackage{lmodern}
\usepackage{amsmath}
\usepackage{booktabs}
\usepackage{array}
\usepackage{tikz}
\usepackage{listings}

\usetikzlibrary{
  arrows.meta,
  positioning,
  shapes.geometric,
  calc
}

% ============================================================
% Colors
% ============================================================

\definecolor{navy}{RGB}{24,43,68}
\definecolor{deepblue}{RGB}{45,88,130}
\definecolor{lightblue}{RGB}{229,238,248}
\definecolor{softgray}{RGB}{245,247,250}
\definecolor{darkgray}{RGB}{38,43,50}
\definecolor{gold}{RGB}{190,140,40}

\setbeamercolor{palette primary}{bg=navy,fg=white}
\setbeamercolor{palette secondary}{bg=deepblue,fg=white}
\setbeamercolor{frametitle}{bg=navy,fg=white}
\setbeamercolor{title}{fg=white}
\setbeamercolor{subtitle}{fg=white}
\setbeamercolor{author}{fg=white}
\setbeamercolor{date}{fg=white}
\setbeamercolor{titlelike}{fg=white}
\setbeamercolor{block title}{bg=deepblue,fg=white}
\setbeamercolor{block body}{bg=lightblue,fg=darkgray}
\setbeamercolor{normal text}{fg=darkgray,bg=white}

% ============================================================
% General formatting
% ============================================================

\setlength{\parskip}{0.35em}

\newcommand{\key}[1]{\textbf{#1}}
\newcommand{\bits}[1]{\texttt{#1}}
\newcommand{\hex}[1]{\texttt{0x#1}}
\newcommand{\code}[1]{\texttt{#1}}
\newcommand{\iron}{\(T_{CPU}=IC\times CPI\times T_{cycle}\)}

% ============================================================
% Metadata
% ============================================================

\title[Topic 2]{
  RISC vs. CISC:\\
  Instruction Set Design Principles \& Trade-offs
}

\subtitle{
  ISA choices, implementation costs,
  and performance consequences
}

\author{
  Computer Systems Architecture -- ARM AArch64
}

\date{}

% ============================================================
% Document
% ============================================================

\begin{document}

% ============================================================
% Slide 1
% ============================================================

\begin{frame}[plain]

\begin{tikzpicture}[remember picture,overlay]

\fill[navy]
  (current page.south west)
  rectangle
  (current page.north east);

\fill[gold]
  (current page.south west)
  rectangle
  ([yshift=.12cm]current page.south east);

\end{tikzpicture}

\vspace{1.0cm}

\begin{minipage}{0.86\textwidth}

{\color{gold}
 \large
 CSCI 6461 \;|\; TOPIC 2
 \par}

\vspace{0.28cm}

{\color{white}
 \fontsize{25}{30}\selectfont
 \bfseries
 RISC vs. CISC:\\[0.08cm]
 Instruction Set Design\\[0.08cm]
 Principles \& Trade-offs
 \par}

\vspace{0.35cm}

{\color{white!78}
 \large
 ISA choices, implementation costs,
 and performance consequences
 \par}

\end{minipage}

\vfill

{\color{white}
 \normalsize
 Dr. J. Burns
 \par}

{\color{white!65}
 \small
 Computer Systems Architecture
 \par}

\vspace{0.45cm}

\end{frame}

% ============================================================
% Slide 2
% ============================================================

\begin{frame}{What Is an Instruction Set Architecture?}

\begin{itemize}
\item An ISA is the programmer-visible specification of a processor.
\item It defines instructions, registers, data types, addressing modes,
      and architectural state.
\item Compilers generate code for the ISA; processors implement the ISA.
\item ISA design therefore sits directly at the hardware/software boundary.
\item A good ISA must remain useful for decades while implementations
      evolve underneath it.
\item RISC and CISC are different historical answers to the question:
      \key{what should the ISA expose?}
\end{itemize}

\end{frame}

% ============================================================
% Slide 3
% ============================================================

\begin{frame}{Why ISA Design Matters}

\begin{columns}[T,onlytextwidth]

\begin{column}{0.46\textwidth}
\begin{block}{Software consequences}
\begin{itemize}
\item Instruction count
\item Code size
\item Compiler complexity
\item Calling conventions
\item Binary compatibility
\end{itemize}
\end{block}
\end{column}

\hfill

\begin{column}{0.46\textwidth}
\begin{block}{Hardware consequences}
\begin{itemize}
\item Decode complexity
\item Datapath organization
\item Pipeline regularity
\item Control logic
\item Energy per instruction
\end{itemize}
\end{block}
\end{column}

\end{columns}

\vspace{0.4em}

\begin{flushleft}
\key{ISA design redistributes complexity between software and hardware
rather than eliminating it.}
\end{flushleft}

\end{frame}

% ============================================================
% Slide 4
% ============================================================

\begin{frame}{The Instruction-Set Design Problem}

\begin{itemize}
\item Should one instruction perform a small primitive operation or a
      larger sequence of work?
\item Should arithmetic instructions access memory directly or operate
      only on registers?
\item Should instructions have one fixed length or many lengths?
\item How many addressing modes should hardware understand?
\item Should frequently used operations receive special instructions
      even if they complicate decode?
\item RISC and CISC emerged from different answers shaped by the
      technology of their eras.
\end{itemize}

\end{frame}

% ============================================================
% Slide 5
% ============================================================

\begin{frame}{Historical Constraints Shaped Early ISAs}

\begin{itemize}
\item Early main memory was expensive, slow, and physically limited.
\item Compilers were less sophisticated than modern optimizing compilers.
\item Hardware designers wanted programs to occupy fewer bytes and make
      fewer memory accesses.
\item Rich instructions could express substantial work in a compact encoding.
\item Microprogrammed control allowed complex instructions to be implemented
      as sequences of simpler internal operations.
\item These pressures strongly encouraged what later became known as the
      CISC style.
\end{itemize}

\end{frame}

% ============================================================
% Slide 6
% ============================================================

\begin{frame}{The Origins of CISC}

\begin{itemize}
\item CISC stands for \key{Complex Instruction Set Computer}.
\item The philosophy favored rich instruction sets with many operations
      and addressing modes.
\item Individual instructions could perform multi-step tasks such as
      memory access plus arithmetic.
\item Variable-length encodings allowed common instructions to be compact
      while retaining expressive instructions.
\item The goal was often to reduce program size and narrow the semantic
      gap between languages and machine code.
\item The x86 ISA is the most prominent surviving example of a CISC lineage.
\end{itemize}

\end{frame}

% ============================================================
% Slide 7
% ============================================================

\begin{frame}{The CISC Design Philosophy}

\begin{itemize}
\item Provide many instructions so software can request specialized
      operations directly.
\item Permit operands to reside in registers or memory in many instruction forms.
\item Support multiple addressing modes to make data structures convenient
      to access.
\item Use variable instruction lengths to encode simple cases compactly
      and complex cases when needed.
\item Allow one architectural instruction to correspond to several
      internal hardware actions.
\item Accept more complicated decoding and control in exchange for
      expressive machine instructions.
\end{itemize}

\end{frame}

% ============================================================
% Slide 8
% ============================================================

\begin{frame}{CISC Example: More Work per Instruction}

\begin{block}{Conceptual example}
Imagine an ISA instruction that computes
\code{a = a + memory[b + 8]}.
\end{block}

\begin{itemize}
\item One architectural instruction identifies an address, reads memory,
      performs arithmetic, and writes a result.
\item The instruction may encode an operation, base register,
      displacement, and destination.
\item Fewer instructions are visible in the program.
\item Internally, hardware may still perform address generation,
      memory access, addition, and write-back separately.
\item The architectural instruction is compact semantically even if its
      implementation is multi-step.
\item This distinction between \key{architectural instruction} and
      \key{internal work} is important.
\end{itemize}

\end{frame}

% ============================================================
% Slide 9
% ============================================================

\begin{frame}{Why Complex Instructions Can Be Attractive}

\begin{itemize}
\item Fewer instructions can reduce instruction-fetch traffic.
\item Rich addressing modes can express common data-access patterns directly.
\item Compact encodings improve code density, which can benefit
      instruction-cache behavior.
\item Some specialized operations map naturally to common software idioms.
\item Backward compatibility can preserve huge software ecosystems over
      many processor generations.
\item The attraction is not complexity for its own sake, but accomplishing
      useful work with fewer architectural instructions or bytes.
\end{itemize}

\end{frame}

% ============================================================
% Slide 10
% ============================================================

\begin{frame}{The Cost of Instruction Complexity}

\begin{itemize}
\item Variable instruction length makes instruction boundaries harder
      to find quickly.
\item Many addressing modes increase decoder and control complexity.
\item Instructions may require very different execution times and
      hardware resources.
\item Irregular operand rules complicate pipelines and dependency handling.
\item Complex decode consumes area and energy before useful execution begins.
\item These costs became increasingly important as processors became
      deeply pipelined.
\end{itemize}

\end{frame}

% ============================================================
% Slide 11
% ============================================================

\begin{frame}{The Emergence of RISC}

\begin{itemize}
\item Researchers observed that compilers used only a subset of many
      complex instruction sets frequently.
\item Simple operations often dominated real compiled workloads.
\item Semiconductor integration made larger register files increasingly practical.
\item Faster pipelines rewarded instructions with predictable execution behavior.
\item Compiler technology improved enough to manage more of the scheduling
      and instruction-selection burden.
\item These observations motivated Reduced Instruction Set Computer designs.
\end{itemize}

\end{frame}

% ============================================================
% Slide 12
% ============================================================

\begin{frame}{The RISC Design Philosophy}

\begin{itemize}
\item Favor regular instructions that are straightforward to fetch,
      decode, and execute.
\item Perform arithmetic primarily on registers.
\item Access memory explicitly through load and store instructions.
\item Use relatively regular instruction encodings.
\item Provide enough registers to keep active values close to the execution units.
\item Let compilers combine simple operations into efficient instruction sequences.
\end{itemize}

\end{frame}

% ============================================================
% Slide 13
% ============================================================

\begin{frame}{Berkeley RISC and Stanford MIPS}

\begin{itemize}
\item RISC research accelerated during the late 1970s and early 1980s.
\item Berkeley RISC demonstrated the value of simple instructions and
      large register sets.
\item Stanford MIPS emphasized pipeline-friendly instruction behavior.
\item Both projects challenged the assumption that increasingly complex
      ISAs were necessary for performance.
\item Their ideas strongly influenced later commercial processor families.
\item ARM emerged in the same broader movement toward simpler,
      regular instruction-set design.
\end{itemize}

\end{frame}

% ============================================================
% Slide 14
% ============================================================

\begin{frame}{RISC Principle: Simple and Regular Instructions}

\begin{itemize}
\item Regularity makes instruction behavior easier for hardware to predict.
\item Similar instruction formats simplify field extraction during decode.
\item Operations tend to have clearly defined sources and destinations.
\item Simpler instruction semantics reduce special-case control paths.
\item Regularity becomes particularly valuable when many instructions
      overlap in a pipeline.
\item The goal is not necessarily fewer instruction types, but predictable structure.
\end{itemize}

\end{frame}

% ============================================================
% Slide 15
% ============================================================

\begin{frame}{RISC Principle: Load/Store Architecture}

\begin{itemize}
\item Arithmetic and logical instructions operate primarily on registers.
\item Loads transfer data from memory into registers.
\item Stores transfer data from registers back to memory.
\item Arithmetic instructions therefore do not normally need to wait
      directly for arbitrary memory operands.
\item Memory behavior becomes concentrated in explicit instruction classes.
\item AArch64 follows this load/store organization.
\end{itemize}

\end{frame}

% ============================================================
% Slide 16
% ============================================================

\begin{frame}{RISC Principle: Register-Based Computation}

\begin{itemize}
\item Registers are much faster to access than main memory.
\item A large register set allows temporary values to remain near execution hardware.
\item Compilers attempt to assign frequently used program values to registers.
\item Register-to-register operations have regular operand access behavior.
\item Reduced memory traffic can improve both performance and energy efficiency.
\item Register allocation therefore becomes an important compiler responsibility.
\end{itemize}

\end{frame}

% ============================================================
% Slide 17
% ============================================================

\begin{frame}{RISC Principle: Regular Instruction Encoding}

\begin{itemize}
\item Fixed-width instructions make instruction boundaries immediately known.
\item Fetch hardware can retrieve several aligned instructions predictably.
\item Decoder logic knows where common fields are likely to appear.
\item Parallel decoding becomes easier when instruction boundaries are regular.
\item AArch64 uses fixed 32-bit instruction encodings.
\item Encoding regularity is one reason RISC designs map naturally onto pipelines.
\end{itemize}

\end{frame}

% ============================================================
% Slide 18
% ============================================================

\begin{frame}{RISC Principle: Compiler-Friendly ISA Design}

\begin{itemize}
\item RISC deliberately exposes relatively primitive operations to software.
\item Compilers decide how those operations should be combined.
\item Register allocation determines which values remain in fast registers.
\item Instruction scheduling can rearrange independent operations to reduce stalls.
\item A regular ISA gives the compiler predictable building blocks.
\item Some complexity therefore moves from instruction semantics into
      compiler optimization.
\end{itemize}

\end{frame}

% ============================================================
% Slide 19
% ============================================================

\begin{frame}[fragile]{Worked Example: Evaluating an Expression in RISC}

Suppose:

\[
c = a + b
\]

and the values are stored in memory.

\begin{block}{Conceptual AArch64-style sequence}
\begin{verbatim}
ldr x1, [x0]
ldr x2, [x0, #8]
add x3, x1, x2
str x3, [x0, #16]
\end{verbatim}
\end{block}

\begin{itemize}
\item Memory access and arithmetic are separate operations.
\item The actual addition occurs between register operands.
\item Four architectural instructions are visible.
\end{itemize}

\end{frame}

% ============================================================
% Slide 20
% ============================================================

\begin{frame}[fragile]{Hypothetical CISC Equivalent}

The same conceptual operation might be exposed as:

\begin{block}{Hypothetical CISC instruction}
\begin{verbatim}
add_mem [c], [a], [b]
\end{verbatim}
\end{block}

Conceptually, this one instruction performs:

\begin{enumerate}
\item Read operand \code{a} from memory.
\item Read operand \code{b} from memory.
\item Add the two values.
\item Write the result to memory location \code{c}.
\end{enumerate}

\begin{block}{Important distinction}
This is a deliberately hypothetical instruction used to illustrate
CISC philosophy; it is not being presented as an actual x86 instruction.
\end{block}

\end{frame}

% ============================================================
% Slide 21
% ============================================================

\begin{frame}{RISC vs. CISC: Same Computation, Different Interface}

\begin{columns}[T,onlytextwidth]

\begin{column}{0.46\textwidth}
\begin{block}{RISC view}
Several explicit operations:
\begin{itemize}
\item load
\item load
\item add
\item store
\end{itemize}
\end{block}
\end{column}

\hfill

\begin{column}{0.46\textwidth}
\begin{block}{CISC view}
One richer architectural operation may:
\begin{itemize}
\item identify memory operands
\item perform address generation
\item perform arithmetic
\item update memory
\end{itemize}
\end{block}
\end{column}

\end{columns}

\vspace{0.3cm}

The CISC sequence may contain fewer architectural instructions, but
that does not imply less internal processor work.

\end{frame}

% ============================================================
% Slide 22
% ============================================================

\begin{frame}{Instruction Count: Is Fewer Always Better?}

\begin{itemize}
\item A complex instruction can replace several simpler instructions.
\item This may reduce dynamic instruction count.
\item Fewer instructions can also reduce instruction-fetch bandwidth.
\item But instruction count is only one term in processor performance.
\item A complex instruction may require more cycles or more decoding work.
\item Therefore lower instruction count does not automatically imply
      lower execution time.
\end{itemize}

\end{frame}

% ============================================================
% Slide 23
% ============================================================

\begin{frame}{Cycles per Instruction: The Other Side}

\begin{itemize}
\item CPI measures average clock cycles consumed per executed instruction.
\item Simple instructions typically require fewer distinct execution resources and exhibit 
more regular latency behavior
\item Complex instructions may require several internal operations.
\item Modern superscalar processors can execute multiple simple operations
      concurrently.
\item Memory delays can dominate the CPI of either ISA style.
\item Instruction count and CPI must therefore be considered together.
\end{itemize}

\end{frame}

% ============================================================
% Slide 24
% ============================================================

\begin{frame}{Clock Cycle Time and Instruction Complexity}

\begin{itemize}
\item Processor clock period is constrained by critical hardware paths.
\item More complicated decode and control can increase front-end delay.
\item Simpler instruction formats can make high-speed decoding easier.
\item Modern processors use pipelining to divide complex work across stages.
\item A longer pipeline can support higher clock frequency but introduces
      other penalties.
\item ISA complexity can therefore influence clock design indirectly
      through implementation choices.
\end{itemize}

\end{frame}

% ============================================================
% Slide 25
% ============================================================

\begin{frame}{The CPU Performance Equation (“Iron Law”)}

\[
\boxed{
T_{CPU}
=
IC
\times
CPI
\times
T_{cycle}
}
\]

\begin{itemize}
\item CISC may reduce \key{instruction count}.
\item RISC regularity may help reduce \key{CPI}.
\item Simpler decode and pipeline structures may help reduce
      \key{clock cycle time}.
\item These effects can oppose one another.
\item No single term determines performance by itself.
\item ISA design must therefore be evaluated using complete workload execution.
\end{itemize}

\end{frame}

% ============================================================
% Slide 26
% ============================================================

\begin{frame}{Worked Performance Comparison}

Assume two implementations execute the same workload.

\begin{center}
\begin{tabular}{lccc}
\toprule
 & IC & CPI & Clock Rate \\
\midrule
RISC & $1.4\times10^9$ & 1.1 & 3.0 GHz \\
CISC & $1.0\times10^9$ & 1.6 & 3.0 GHz \\
\bottomrule
\end{tabular}
\end{center}

\[
T_{\text{RISC}}
=
\frac{1.4\times10^9 \times 1.1}{3.0\times10^9}
\approx 0.513\text{ s}
\]

\[
T_{\text{CISC}}
=
\frac{1.0\times10^9 \times 1.6}{3.0\times10^9}
\approx 0.533\text{ s}
\]

The RISC machine executes more instructions yet completes the workload sooner.

\end{frame}

% ============================================================
% Slide 27
% ============================================================

\begin{frame}{Code Density and Memory Footprint}

\begin{itemize}
\item Program size depends on both instruction count and bytes per instruction.
\item Variable-length CISC encodings can represent common operations compactly.
\item Fixed-width RISC instructions may require more bytes for some sequences.
\item Higher code density can improve instruction-cache utilization.
\item Better instruction-cache behavior can reduce memory traffic and energy.
\item Code density therefore remains a genuine architectural advantage,
      especially when memory capacity or bandwidth is constrained.
\end{itemize}

\end{frame}

% ============================================================
% Slide 28
% ============================================================

\begin{frame}{Instruction Fetch and Decode Complexity}

\begin{columns}[T,onlytextwidth]

\begin{column}{0.46\textwidth}
\begin{block}{Regular fixed-width encoding}
\begin{itemize}
\item boundaries known immediately
\item easier parallel fetch
\item regular field extraction
\item simpler alignment
\end{itemize}
\end{block}
\end{column}

\hfill

\begin{column}{0.46\textwidth}
\begin{block}{Variable-length encoding}
\begin{itemize}
\item boundaries must be discovered
\item instructions may cross fetch boundaries
\item decode may require more stages
\item more front-end hardware may be required
\end{itemize}
\end{block}
\end{column}

\end{columns}

\vspace{0.3cm}

Modern CISC processors invest substantial hardware in making this
front end fast enough to feed wide execution engines.

\end{frame}

% ============================================================
% Slide 29
% ============================================================

\begin{frame}{ISA Design and Pipelining}

\begin{itemize}
\item Pipelines work best when instruction behavior can be divided into
      regular stages.
\item RISC load/store organization separates memory operations from
      arithmetic operations.
\item Fixed instruction width simplifies fetch and decode stage boundaries.
\item Similar operand structures make hazard detection easier to organize.
\item Irregular instructions can still be pipelined, but may require
      additional translation or control.
\item RISC historically aligned well with the needs of deeply pipelined processors.
\end{itemize}

\end{frame}

% ============================================================
% Slide 30
% ============================================================

\begin{frame}{ISA Design and Compiler Optimization}

\begin{itemize}
\item Compilers translate high-level program behavior into ISA operations.
\item A regular instruction set provides predictable optimization targets.
\item Register allocation determines which program values remain in registers.
\item Instruction scheduling attempts to expose independent operations.
\item Rich CISC instructions can sometimes express useful operations directly.
\item ISA quality therefore includes how effectively real compiler
      toolchains can exploit its features.
\end{itemize}

\end{frame}

% ============================================================
% Slide 31
% ============================================================

\begin{frame}{RISC vs. CISC: Energy and Thermal Trade-offs}

\begin{columns}[T,onlytextwidth]

\begin{column}{0.46\textwidth}
\begin{block}{RISC tendency}
\begin{itemize}
\item Regular encodings can simplify fetch, alignment, and decode.
\item Simpler architectural operations can reduce front-end work per instruction.
\item Explicit load/store behavior can make data movement more predictable.
\end{itemize}
\end{block}
\end{column}

\hfill

\begin{column}{0.46\textwidth}
\begin{block}{CISC tendency}
\begin{itemize}
\item Variable-length or richer instructions can require more complex decode.
\item One instruction may perform several operations or memory accesses.
\item Better code density can reduce instruction-fetch traffic.
\end{itemize}
\end{block}
\end{column}

\end{columns}

\vspace{0.2cm}

\begin{block}{Architectural implication}
ISA style influences energy, but implementation and workload determine
total efficiency.
\end{block}

\end{frame}

% ============================================================
% Slide 32
% ============================================================

\begin{frame}{Energy per Instruction Is Not Energy per Program}

\begin{block}{A useful first-order model}
\[
E_{\text{program}}
\approx
IC \times E_{\text{instruction}}
\]
\end{block}

\begin{itemize}
\item A simpler instruction may consume less energy when fetched,
      decoded, and executed.
\item But a RISC sequence may require more architectural instructions
      to complete the same task.
\item A denser CISC sequence may reduce instruction-cache accesses.
\item Memory behavior can dominate energy when cache misses reach DRAM.
\item The relevant metric is total energy for the workload.
\item This mirrors the Iron Law: optimizing one factor alone can produce
      the wrong conclusion.
\end{itemize}

\end{frame}

% ============================================================
% Slide 33
% ============================================================

\begin{frame}{Power Becomes Heat --- and Heat Requires Cooling}

\begin{columns}[T,onlytextwidth]

\begin{column}{0.44\textwidth}
\begin{block}{Energy and power}
\[
P = \frac{E}{t}
\]

Higher switching activity, voltage, frequency, and active hardware
structures increase power demand.
\end{block}
\end{column}

\hfill

\begin{column}{0.47\textwidth}
\begin{block}{System consequence}
\centering
architectural complexity\\
$\downarrow$\\
transistor switching\\
$\downarrow$\\
power consumption\\
$\downarrow$\\
heat generation\\
$\downarrow$\\
cooling requirement
\end{block}
\end{column}

\end{columns}

\vspace{0.15cm}

\begin{itemize}
\item Cooling ranges from passive heat spreaders to fans, liquid cooling,
      and data-center HVAC.
\item Thermal limits can force frequency reduction, so power can ultimately
      constrain performance.
\end{itemize}

\end{frame}

% ============================================================
% Slide 34
% ============================================================

\begin{frame}{Performance per Watt: The More Useful Question}

\begin{itemize}
\item Processor design increasingly optimizes useful work per unit of energy.
\item Simpler decode and regular execution can contribute to efficiency.
\item Code density, caches, speculation, vector units, and memory traffic
      can reverse simple ISA-level expectations.
\item Process technology, voltage, frequency, and microarchitecture often
      matter more than the RISC/CISC label.
\item ARM's success now spans phones, laptops, and cloud servers.
\item Therefore, ``ARM is efficient because RISC is simpler'' is an
      incomplete architectural explanation.
\end{itemize}

\begin{block}{Architectural implication}
Compare \key{energy per workload}, \key{performance per watt}, and
\key{thermal constraints}; ISA labels are only part of the evidence.
\end{block}

\end{frame}

% ============================================================
% Slide 35
% ============================================================

\begin{frame}{Why Modern x86 Looks RISC-Like Inside}

\begin{itemize}
\item Modern x86 cores preserve the complex x86 ISA for software compatibility.
\item The front end decodes many x86 instructions into simpler internal
      micro-operations (\(\mu\)ops).
\item Those \(\mu\)ops are scheduled onto regular execution resources.
\item This separates a complex architectural interface from a more regular
      internal execution model.
\item Some x86 instructions decode into one \(\mu\)op; others require several.
\item ISA style and microarchitecture style are therefore not the same thing.
\end{itemize}

\end{frame}

% ============================================================
% Slide 36
% ============================================================

\begin{frame}{Why Modern RISC ISAs Are Not Necessarily ``Simple''}

\begin{itemize}
\item Modern RISC ISAs include SIMD, cryptography, virtualization,
      atomics, and system-control instructions.
\item AArch64 contains many instruction classes and sophisticated
      addressing features.
\item High-performance ARM cores use branch prediction, out-of-order
      execution, speculation, and large caches.
\item None of this violates the central RISC ideas of regular encoding
      and load/store organization.
\item ``Reduced'' should be understood historically as architectural
      regularity, not a tiny instruction count.
\item Calling a modern ARM core ``simple'' confuses ISA philosophy with
      implementation sophistication.
\end{itemize}

\end{frame}

% ============================================================
% Slide 37
% ============================================================

\begin{frame}{RISC Beyond Phones: Laptops}

\begin{itemize}
\item Modern RISC is no longer confined to embedded controllers or phones.
\item Apple Silicon brought Arm-based processors into mainstream laptop
      and desktop computing.
\item Current MacBook Pro systems use Apple M5-family processors.
\item These systems target high performance while retaining strong
      energy efficiency and battery life.
\item Laptops therefore demonstrate the importance of performance per watt,
      not merely low power.
\item RISC principles can scale from constrained devices to high-performance
      personal computers.
\end{itemize}

\end{frame}

% ============================================================
% Slide 38
% ============================================================

\begin{frame}{RISC in the Data Center: Google Axion}

\begin{itemize}
\item Cloud providers now deploy Arm-based processors for general-purpose
      server workloads.
\item Google Axion is a custom Arm-based CPU designed for Google Cloud
      data-center workloads.
\item Target workloads include web servers, databases, analytics,
      media processing, and CPU-based AI.
\item Google has described deployment for Bigtable, Spanner, BigQuery,
      Blobstore, Pub/Sub, Earth Engine, and YouTube Ads.
\item These are latency-sensitive and throughput-oriented production workloads.
\item At data-center scale, performance per watt affects server density,
      electricity cost, heat generation, and cooling infrastructure.
\end{itemize}

\end{frame}

% ============================================================
% Slide 39
% ============================================================

\begin{frame}{Data-Center Workloads: What Runs on Arm?}

\begin{columns}[T,onlytextwidth]

\begin{column}{0.46\textwidth}

\begin{block}{Storage and databases}
\begin{itemize}
\item Spanner
\item Bigtable
\item Cloud SQL / AlloyDB-class workloads
\end{itemize}
\end{block}

\vspace{0.2cm}

\begin{block}{Analytics and pipelines}
\begin{itemize}
\item BigQuery
\item Dataflow-style pipelines
\item Batch and containerized jobs
\end{itemize}
\end{block}

\end{column}

\hfill

\begin{column}{0.46\textwidth}

\begin{block}{Serving and platforms}
\begin{itemize}
\item YouTube Ads
\item Web and application servers
\item Containerized microservices
\end{itemize}
\end{block}

\vspace{0.2cm}

\begin{block}{AI and inference}
\begin{itemize}
\item CPU-based inference
\item Preprocessing and data movement
\item Non-accelerator-dominated services
\end{itemize}
\end{block}

\end{column}

\end{columns}

\end{frame}

% ============================================================
% Slide 40
% ============================================================

\begin{frame}{ARM AArch64 as a Modern RISC ISA}

\begin{itemize}
\item AArch64 uses fixed 32-bit instruction encodings.
\item Most arithmetic and logical operations use register operands.
\item Memory is accessed primarily through explicit load and store instructions.
\item Thirty-one general-purpose integer registers support
      register-oriented computation.
\item Instruction formats are structured so common fields remain
      relatively regular.
\item These properties make AArch64 useful for studying the connection
      between ISA and datapath design.
\end{itemize}

\begin{block}{Next topic}
The course now moves from ISA design philosophy to the actual AArch64
programmer's model and assembly language.
\end{block}

\end{frame}

% ============================================================
% Slide 41
% ============================================================

\begin{frame}{RISC vs. CISC: What Actually Matters Today?}

\begin{itemize}
\item The RISC/CISC labels are less predictive than they were in the 1980s.
\item Modern CISC processors use sophisticated internal decomposition
      and execution engines.
\item Modern RISC processors contain complex microarchitectures and rich
      ISA extensions.
\item Compatibility, code density, compiler ecosystem, implementation
      quality, and workload matter enormously.
\item The enduring RISC contribution is emphasis on regularity,
      load/store organization, and compiler-visible simple operations.
\item Architectural choices should be evaluated quantitatively rather
      than through slogans.
\end{itemize}

\end{frame}

% ============================================================
% Slide 42
% ============================================================

\begin{frame}{Synthesis: ISA Choices Are Trade-offs}

\begin{itemize}
\item CISC tends to expose richer operations and often achieves strong
      code density.
\item RISC tends to expose more regular operations that are easier to
      decode and pipeline.
\item Neither lower instruction count nor higher clock rate guarantees
      better performance.
\item The Iron Law requires instruction count, CPI, and cycle time to be
      considered together.
\item Energy, heat, cooling, and performance per watt extend the analysis
      beyond execution time alone.
\item AArch64 provides a modern RISC case study for the remainder of the course.
\end{itemize}

\begin{block}{Design question}
Do not ask ``Which philosophy is faster?''

Ask: \key{What trade-off does this ISA choice create, and what does the
measured workload tell us?}
\end{block}

\end{frame}

\end{document}